\documentclass{article}
\usepackage[utf8]{inputenc}
\usepackage[T1]{fontenc}
\usepackage[margin=2cm]{geometry}
\usepackage{hyperref}
\sloppy
\begin{document}
\section*{Historical source}
Encyclopaedia Britannica scan, source row 2051. Parse 3.5 extracted the text. Super 3.5 VL corrected it from an enlarged crop of each block. It can still contain recognition errors.
Page headers, footers, and page numbers are removed.

\subsection*{Untitled continuation (E0178)}
\textit{Continues an entry from an earlier page.}

LOM [ 168 ]

Lombards, judging that a favourable opportunity now offered of making himself king of Italy by her means, imparted his design to her, and declared his intention to marry her, provided, by some means or other, the despatches Helmchild.—Rosamund, highly pleased with the proposal, resolved to satisfy her ambition by getting rid of the person whom he had married in order to gratify her revenge. Accordingly, having prepared a strong poison, he mixed it with wine, and gave it to her husband as he came out of the bath, and called for drink, according to his custom. Helmchild had not half emptied the cup, when, by the sudden and strange operation which he felt in his bowels, he concluded what it was; and, with his sword pointed at the queen's breast, compelled her to drink the rest. The poison had the same effect on both; for they died in a few hours. Longinus, on the death of the queen, laid aside all thoughts of making himself king of Italy, and sent the king's treasure to Constantinople, together with Albifonda, the daughter of Alboinus by Rosamund, whom he had brought alone with her.

After the death of Alboinus, the Lombards chose Clephis, one of the nobility, for their king. He was murdered after a short reign of 18 months; upon which ensued an interregnum of 10 years, as related under the article ITALY, N° 32. During this time, they extended their conquests in that country; but at last the Romans, jealous of their progress, resolved to put a stop to their victories, and, if possible, to drive them quite out. For this purpose, they designed not only to employ their own force, but entered into alliance with the Franks; which so alarmed the Lombards that they re-established the monarchical form of government among themselves, and chose Autharis the son of Clephis for their king. This monarch, considering that the power of the dukes, who had governed Lombardy for the space of 10 years, was during that length of time very much established, and that they would not probably be willing to part with the authority which they had so long enjoyed, allowed them to continue in their government; but obliged them to contribute one moiety of their revenues towards the maintenance and support of his royal dignity, suffering them to dispose of the other as they thought proper. He reserved to himself the supreme dominion and authority; and took an oath of the dukes, that in time of war, they would readily assist him to the utmost of their power. Though he could remove the dukes at pleasure, yet he deprived none of them of their dukedoms, except in cases of treason; nor gave them to others, except when their male issue sailed. Having settled matters in this manner with the dukes, he enacted several wholesome laws against theft, rape, murder, adultery, and other vices which prevailed among his subjects, and was the first of the Lombard kings who embraced Christianity. Most of his subjects followed the example of their monarch; but as they were all instructed by Arian bishops, they continued long infected with that heresy; which occasioned great disputes between them and the orthodox bishops of the cities subject to them.

Written laws introduced.

laws among these people. Before his time they had been governed only by tradition: but Rotharis, in imitation of the Romans and Goths, undertook the publishing of written laws; and to those which he enacted, many were added by the succeeding princes. Grotius prefers the method which the Lombards followed in making laws, to that which was professed by the Romans themselves. Among the latter the emperor was the sole lawgiver; so that whatever pleased him had the force of a law. But the Lombard kings did not assume that power to themselves, since their laws were enacted in public assemblies, convened for that purpose, after they had been maturely examined and approved of by all the lords of the kingdom. From these assemblies were excluded the ecclesiastical order, and the people: so that the legislative power was lodged in the king and nobles alone.

The reign of Rotharis is remarkable, not only for his introducing written laws among his subjects, but for the conquests he made, and the successful wars carried on with the exarch of Ravenna, whom he totally defeated in several engagements, and made himself master of some part of his territories. This monarch died in 652; and the affairs of the Lombards went on prosperously, till the ambition of Luitprand laid the foundation of the total ruin of his kingdom. He ascended the throne of Lombardy in 711, and Luitprand's watched all opportunities of enlarging his dominions at the expense of the emperor's. Of this, a fair opportunity offered in 716: for the emperor Leo Isauricus, who at that time reigned in the east, having, by his famous edict, forbidden the worship of images, and ordered them to be everywhere pulled down, the people were so provoked at that innovation, that, in several places, they openly revolted, and, falling upon the emperor's officers, drove them out of the cities. In the east, Germanus, patriarch of Constantinople, opposed the emperor's design with great warmth; but Leo caused him to be deposed, and Anastasius to be raised to that see in his room, ordering at the same time all the images in the imperial city to be pulled down and publicly burnt. He strictly enjoined his officers in the west, especially the exarch of Ravenna, to see his edict punctually obeyed in their respective governments. In compliance with these orders, Scholasticus, then exarch, began to pull down the images in all the churches and public places in Ravenna; which incensed the superstitious multitude to such a degree, that taking arms, they openly declared they would rather renounce their allegiance to the emperor than the worship of images.

Thus a kind of civil war being kindled in the city, Luitprand thought he had now a favourable opportunity of making himself master of the feast of the exarch, not doubting but the conquest of such an important place would be followed by that of the whole exarchate. Having therefore drawn together all his forces, he unexpectedly appeared before Ravenna, and closely besieged it. The exarch little expected such a surprise, as a friendly correspondence had been maintained for many years between the exarchs and the Lombard kings. However, he defended the place with such courage and resolution, that Luitprand, despairing of success, broke up the siege and led his army against Classis, at a small distance from Ravenna, which


\section*{Source record}
\begin{description}
\item[Scan ID] \texttt{britannica:309f765b1ecc83e0a7f7762e44b4a927a7555443:train:2051}
\item[Image file] \texttt{page-2051-588566464df4.jpg}
\item[Image SHA-256] \texttt{588566464df4b426ffc48c73095181f8040fb4c5dfb1bf3ef4738390589d2302}
\item[Dataset] NationalLibraryOfScotland/encyclopaedia\_britannica\_illustrated at revision \texttt{309f765b1ecc83e0a7f7762e44b4a927a7555443} (CC0-1.0)
\item[Extraction] model \texttt{nvidia/nemotron-3.5-parse}, raw reply \texttt{extract/raw/page-2051-588566464df4.json}
\item[Raw reply SHA-256] \texttt{556558e35b6f5b9b7c9a9d70a2e4bab1c0616de2b22c74409d60b4e2e84dbbbf}
\item[Correction] model \texttt{nvidia/Nemotron-3.5-Super-VL}, block crops, raw replies \texttt{correct/raw/crop\_nothink/page-2051-588566464df4-b*.json}; super\_vl: 8
\item[Running header] none
\end{description}
\end{document}
